\section*{الفصل \ref{ch:b3:microcanonical-ensemble} --- الجمهور الميكروقانوني}

\begin{solution}{exo:b3:microcanonical-ensemble:1}
(أ) $1, 4, 6, 4, 1$. (ب) على $16$: $6.25\%$، $25\%$، $37.5\%$،
$25\%$، $6.25\%$. (ج) $S/k_{\text{B}} = 0$، $\ln4$، $\ln6$، $\ln4$،
$0$. (د) $n = 2$، ولا يحمل إلا $37.5\%$ من الاحتمال: فعند
$N = 4$ لا يكون ``التوازن'' إلا أغلبية نسبية --- والحدّة
التي تشرّع الترموديناميك تُشترى بفضل $N$ الكبيرة.
\end{solution}

\begin{solution}{exo:b3:microcanonical-ensemble:2}
(أ) $N = 10$: $13.0$ مقابل $15.1$ (بفارق $14\%$)؛ و $N = 100$:
$360.5$ مقابل $363.7$ (بفارق $0.9\%$). (ب) الحدّ المهمَل هو
$\tfrac12\ln(2\pi N)$: أي خطأ نسبي $\sim\ln N/N$. (ج) طبّق
ستيرلنغ على العواملات الثلاثة. (د) عند $N = 10^{23}$ يكون الحدّ
المطروح $\sim26$ مقابل $N\ln2 \sim 10^{23}$: أي اثنتان وعشرون رتبة
تحت الحدّ الرئيسي.
\end{solution}

\begin{solution}{exo:b3:microcanonical-ensemble:3}
(أ) $\ln(\Omega_1\Omega_2) = \ln\Omega_1 + \ln\Omega_2$. (ب)
$\Delta S = k_{\text{B}}\ln2$ --- أي جسيم واحد عُرض عليه حجم
مضاعف. (ج) $S = 156\,k_{\text{B}} = \qty{2.2e-21}{J/K}$. (د)
تحمل الإنتروبيات الترموديناميكية عوامل مقدارها $10^{23}$: فكل
الخلط في كل ملهى قمار على الأرض لا يُرى إلى جانب نفَس
من هواء دافئ.
\end{solution}

\begin{solution}{exo:b3:microcanonical-ensemble:4}
(أ) $\Delta S = N_{\text{A}}k_{\text{B}}\ln2 = R\ln2 =
\qty{5.76}{J/K}$. (ب) $2^{-N_{\text{A}}} \sim 10^{-1.8\times
10^{23}}$. (ج) حتى لو عوينت التشكيلات كل \qty{e-10}{s}،
يقزّم زمن الانتظار المتوقَّع المقدارَ \qty{4e17}{s} بعامل
\emph{أُسّه} من ثلاثة وعشرين رقمًا. (د) ليس مستحيلًا ---
بل غير قابل للمشاهدة: فالقانون احتمالي من حيث المبدأ ومطلق
في أيّ عالم يحوي راصدين ذوي صبر منتهٍ.
\end{solution}

\begin{solution}{exo:b3:microcanonical-ensemble:5}
(أ) خذ لوغاريتم الجداء. (ب) بمضاعفة $V, N$ عند
$E/N$ ثابت: من دون $N!$ تكسب $S$ فائضًا زائدًا من نمط $N\ln2$
--- وهو إخفاق غيبس. (ج) ستيرلنغ على $\ln N!$ يخرج
التركيبين $V/N$ و $E/N$. (د) خلية بلانك $h^3$ تثبّت
صفر الإنتروبيا؛ وتطابق الجسيمات يثبّت $N!$ --- أي معطيان
كموميان يختبئان في الترموديناميك الكلاسيكية.
\end{solution}

\begin{solution}{exo:b3:microcanonical-ensemble:6}
(أ) $1/T = \tfrac32 Nk_{\text{B}}/E$. (ب) $P/T =
Nk_{\text{B}}/V$: أي قانون الغاز المثالي، من الإحصاء. (ج) عند
$S$ ثابت: يثبت $\ln(V/N) + \tfrac32\ln(E/N)$، ومنه $VE^{3/2}$ ثابت؛
ومع $E \propto T$: يثبت $TV^{2/3}$، وهو يكافئ
$PV^{5/3}$. (د) قانون الغاز، وعلاقة الطاقة بدرجة الحرارة،
والأُسّ الكظوم --- وهي ثلاث ركائز تجريبية في كتاب السنة
الأولى --- صارت الآن مبرهنات.
\end{solution}

\begin{solution}{exo:b3:microcanonical-ensemble:7}
(أ) التناظر (كتلتان متساويتان). (ب) لكلٍّ $\ln\Omega \propto
\tfrac32 N\ln E_i$: فينتج نشرُه حول التوزيع المتساوي
غاوسيةً في $x$ تباينها $\sim1/N$. (ج) $x_{\text{جذري}} \sim
\num{e-11}$. (د) $\Delta T/T \sim x_{\text{جذري}}$: أي
نانو-نانوكلفنات --- فالتذبذبات موجودة، وهي غير قابلة للقياس
البتة على مقياس المخبر.
\end{solution}

\begin{solution}{exo:b3:microcanonical-ensemble:8}
(أ) $\Omega = \binom Nn$، و $S = k_{\text{B}}\ln\binom Nn$. (ب)
$1/T = \partial S/\partial E = (k_{\text{B}}/\epsilon)[\ln(N-n) -
\ln n]$؛ ثم حلّ من أجل $n/N$. (ج) ترتفع $E(T)$ من $0$ إلى قيمة
الإشباع $N\epsilon/2$؛ وتنعدم $C = \dd E/\dd T$ عند $T$ المنخفضة
(فلا يُستطاع دفع ثمن كمّ) وعند $T$ المرتفعة (فالمستويان مملوءان
بالسواء، ولم يبقَ ما يُمتصّ): وبينهما حدبة --- أي البصمة
المسعّرية لأيّ جمهرة ذات مستويين، تُستعمل لكشف الشوائب
المغناطيسية في الأجسام الصلبة. (د) إنها إشغال حراري ذو مستويين
--- ومع $\epsilon \to E - \mu$ ستعود حرفيًا في صورة
توزيع فيرمي--ديراك.
\end{solution}

\begin{solution}{exo:b3:microcanonical-ensemble:9}
(أ) $1/T \propto \ln[(N-n)/n] < 0$ بالضبط حين $n > N/2$. (ب)
$T = \epsilon/[k_{\text{B}}\ln(2/3)] = -\qty{2.9}{K}$. (ج) الجسم
ذو $T$ السالبة \emph{يكسب} إنتروبيا بفقد الطاقة، والجسم
ذو $T$ الموجبة يكسبها بتلقّيها: فيرتفع الإحصاءان معًا حين تسيل
الطاقة من السالب $\to$ الموجب --- أي إن الجملة المقلوبة هي
المتبرّع الشامل، ومن ثمّ أسخن من كل شيء. (د) مع طاقة غير محدودة
لا ينقلب $\Omega(E)$ أبدًا: فيبقى $\partial S/\partial E$
موجبًا عند أيّ طاقة منتهية.
\end{solution}

\begin{solution}{exo:b3:microcanonical-ensemble:10}
(أ) يضاعف كل غاز حجمه المتاح: $\Delta S =
2Nk_{\text{B}}\ln2$. (ب) تمتصّ عواملات $N!$ للجمهرتين الممزوجتين
الكسبَ الظاهري بالضبط: $\Delta S = 0$. (ج) زيادة الإنتروبيا هي
\emph{فقدان الفرز}: فلم يتحرك أيّ متغير عياني، لكن إعادة
الفصل تتطلب الآن شغلًا --- أي إن الزيادة مخزونة فاتورةً آجلة. (د)
$W_{\min} = T\Delta S =
2Nk_{\text{B}}T\ln2$؛ ومعامل فصل النظائر (كسلاسل نابذات
اليورانيوم واستخلاص الهليوم-3) يدفع هذا الحدّ الترموديناميكي
الأدنى أضعافًا مضاعفة في الواقع.
\end{solution}

\begin{solution}{exo:b3:microcanonical-ensemble:11}
(أ) $\lambda_T = \qty{1.6e-11}{m}$. (ب) $V/N = k_{\text{B}}T/P =
\qty{4.1e-26}{m^3}$: ومنه $V/N\lambda_T^3 = \num{1.0e7}$، أي
$S/Nk_{\text{B}} = \ln(\num{1.0e7}) + 2.5 = 18.6$. (ج) مطابق تمامًا
للقيمة المسعّرية $18.6$. (د) الإنتروبيا المطلقة تحوي $h$
(عبر $\lambda_T$) وتحوي $N!$: فقياسات الحرارة من
عصر البخار تؤكد، بثلاثة أرقام، كمّ \omterm{def:b3:lagrangian-mechanics:action}{الفعل}
وتطابق الذرات.
\end{solution}

\begin{solution}{exo:b3:microcanonical-ensemble:12}
(أ) $v = (\qty{550}{nm})^3 = \qty{1.7e-19}{m^3}$: ومنه $\bar n =
\num{4.2e6}$، و $\delta_{\text{جذري}} = 1/\sqrt{\bar n} \approx
\num{5e-4}$. (ب) الوسط المنتظم انتظامًا تامًّا لا ينتثر شيئًا
جانبيًا (\cref{pb:b3:perturbation-theory:1}، الجزء الثالث): فهذه
التموّجات $\sqrt N$ التي لا تُختزل في كل خلية ضوئية هي بالضبط
الاضطراب الذي يتيح للسماء أن توجد. (ج) عند $\bar n = 100$: يكون $v
\approx (\qty{16}{nm})^3$. (د) السماء زرقاء لأن الهواء مصنوع
من جزيئات معدودة تتذبذب أعدادها وفق $\sqrt N$ ---
واستعمل أينشتاين هذا بالضبط ليحتجّ بأن الجزيئات حقيقية.
\end{solution}

\begin{solution}{pb:b3:microcanonical-ensemble:1}
\textbf{1.} $n_\pm = \tfrac12(N \pm x/b)$؛ و $\Omega =
N!/n_+!\,n_-!$.
\textbf{2.} عند $x = 0$: تتكبّب السلسلة الحرة؛ وللامتداد الكامل
$\Omega = 1$.
\textbf{3.} ستيرلنغ على ثنائي الحدّ، منشورًا حتى المرتبة الثانية في
$x/Nb$: فتنتج الإنتروبيا الغاوسية المذكورة.
\textbf{4.} ما دامت $E$ مستقلة عن $x$، فلا يمكن لأيّ جذب نحو $x$
الصغيرة أن يأتي إلا من الإحصاء: أي إنتروبي بالبناء.
\textbf{5.} الكبّة $\sim b\sqrt N = \qty{50}{nm}$ مقابل $Nb =
\qty{5}{\micro m}$: فالسلسلة الحرة أقصر مئة مرة
من محيطها --- أي متكرمشة كلها تقريبًا.
\textbf{6.} $N$، أي عدد الوصلات بين وصلات التشابك:
فالفلكنة تقصّر السلاسل الفعّالة، وتصلّب
الشبكة.
\textbf{7.} $f = -T\,\partial S/\partial x =
k_{\text{B}}Tx/Nb^2$.
\textbf{8.} $k = k_{\text{B}}T/Nb^2 = \num{4.14e-21}/(\num{e4}
\times \num{2.5e-19}) = \qty{1.7e-6}{N/m}$ لكل سلسلة.
\textbf{9.} $k \propto T$: فحين يُسخَّن تحت حمل ثابت، يتصلّب الرباط
و\emph{يتقلص}، فيرفع الحمل؛ أما النابض الفولاذي فيليّن
ويترهّل قليلًا لا غير.
\textbf{10.} $f(Nb/2) = k_{\text{B}}T/2b \approx \qty{4}{pN}$؛
ويخفق القانون الخطي حين تقترب $f$ من $k_{\text{B}}T/b \approx
\qty{8}{pN}$، حيث تدنو السلسلة من امتدادها الكامل.
\textbf{11.} $\sim\num{e14}$ سلسلة على التوازي و
$\sim\num{e5}$ طول كبّة على التسلسل على امتداد سنتيمتر:
$k_{\text{رباط}} \sim k_{\text{سلسلة}} \times \num{e14}/\num{e5}
\sim \num{e2}$--$\num{e3}\,\unit{N/m}$ --- أي الملمس المألوف
للرباط المطاطي.
\textbf{12.} صلابة المطاط \emph{هي} $k_{\text{B}}T$ لكل
تشكيلة: فمزيد من الاهتياج يعني مزيدًا من الارتداد؛ أما صلابة الفلز
فطاقة روابط، يرخيها الاهتياج اللاتوافقي.
\textbf{13.} المدّ يقصّ الإحصاء التشكيلي؛ وإذا جرى سريعًا
(بلا تبادل إنتروبيا مع الخارج)، وجب أن تعود الإنتروبيا التشكيلية
المفقودة إنتروبيا حرارية: فيدفأ الرباط ---
على مقياس \qty{1}{K}، تكشفه الشفة.
\textbf{14.} والإطلاق يعيد التشكيلات؛ فيردّ الحساب الحراري
الفرق: فيبرد.
\textbf{15.} الثقل \emph{يرتفع}: فالتقلص عند التسخين هو
التوقيع الإنتروبي (وأما المواد المرنة بالطاقة فتتمدد).
\textbf{16.} النابض الفولاذي يستطيل قليلًا (بالتمدد
الحراري) ويهبط معامله: أي إشارة معاكسة --- فهي مرونة
طاقة لا مرونة إحصاء.
\textbf{17.} الأضلاع المسخَّنة تتقلص، فتجذب كتلة الإطار
عن المركز نحو الجهة الباردة؛ فتعزم الثقالة العجلة المختلّة؛
ويرتخي كل ضلع من جديد وهو يبتعد بالدوران: أي محرك حراري
مادة عمله الإنتروبيا نفسها.
\textbf{18.} $k_{\text{B}}T/b = \num{4.14e-21}/\num{e-7} =
\qty{4e-14}{N} = \qty{0.04}{pN}$: أي بيكونيوتنات وما دونها ---
وهو بالضبط مجال الملاقط الضوئية، فأمكن تتبّع قانون القوة--الامتداد
لجزيء حمض نوويّ واحد نقطةً نقطة (ومطابقته،
مع تدقيقات البند 23).
\textbf{19.} من الإحصاء وحده لا غير: فعند الامتداد $x$ تكون
الحالات المجهرية أقلّ منها عند $x - \dd x$، ويسوق الاهتياج الحراري
السلسلةَ نحو الإحصاء الأكبر؛ و``القوة'' هي
$T$ مضروبة في تدرّج الإنتروبيا.
\textbf{20.} عند $T = 0$ لا شيء يستكشف التشكيلات: فتنعدم
القوة الإنتروبية بانعدام الاهتياج الذي يغذّيها ---
ولكانت المرونة طاقيةً محضة، ولتصرّف المطاط
تصرّف خيط مسترخٍ.
\textbf{21.} في الغاز: $P = T\,\partial S/\partial V$ مع $S \ni
Nk_{\text{B}}\ln V$؛ وفي السلسلة: $f = -T\,\partial S/\partial x$ مع
الإحصاء الغاوسي --- وكلاهما إحصاء مدفوع في وجه
\omterm{def:b3:lagrangian-mechanics:coordinates}{قيد}، ومسعَّر بالمقدار $T$.
\textbf{22.} حرارة المجفّف تدخل الرباط؛ ويتحول جزء منها
شغلًا في التقلص متساوي الحرارة، وتوازن الدفاتر الإنتروبيا
المحمولة داخلًا وخارجًا --- أي محرك حراري مصغَّر، والمبدأ الأول
سليم.
\textbf{23.} الامتدادية المنتهية (فالقوة الحقيقية تتباعد قرب
المدّ الكامل --- وهو تدقيق السلسلة الشبيهة بالدودة) و
التبلور المستحثّ بالإجهاد (فالمحاذاة ترتّب السلاسل،
فتعطي تخلّفًا وإطلاق حرارة يتجاوزان النموذج المثالي).
\textbf{24.} $k_{\text{B}}T/\unit{nm} = \qty{4.1}{pN}$: أي
القوة التي تتبادل عندها الطاقة الحرارية والانزياحات النانومترية
بالتساوي --- فالمحركات الجزيئية والبروتينات المنطوية والحمض النووي
الممدود تعمل كلها في حدود رتبة واحدة منها.
\textbf{25.} عشرة آلاف وصلة عمياء، محصاةً: فينتج \omterm{thm:b3:continuum-elasticity:hooke}{قانون هوك} مع
$k = k_{\text{B}}T/Nb^2$، ودفءٌ عند المدّ، وأثقالٌ ترفعها
الحرارة، وجزيئات مفردة تطيعه عند \qty{0.04}{pN} ---
أي إنتروبيا تعمل قوةً.
\end{solution}
